% TOP_PROBLEM: {"id": 30006694, "title": "The Exponential Time Hypothesis for 3-SAT", "category_id": 15, "category": {"id": 15, "name": "computer_science", "display_name": "Computer Science", "description": "Computational complexity, algorithms, and theoretical CS.", "slug": "computer-science", "order_index": 15, "created_at": "2026-07-31T15:26:25.671Z"}, "status": "open"}
% FIRST_POSTED: 2026-10-26
% ATTEMPT_STATUS: unresolved
% SUBMISSION: {"kind": "research_attempt", "category": "top_problems", "problem_id": "30006694", "model": {"provider": "OpenAI", "version": "GPT 6 Astra Pro (exact release/build unknown)", "reasoning": "Unknown", "generated_on": "2026-10-26"}, "other_models": "No second model reported for exposition. Codex used for repository formatting; exact underlying release/build not recorded.", "tools": "Source inspection and TeX preparation reported in supplied manuscript; repository validation, LuaLaTeX and browser rendering during integration. No Lean execution.", "target": "Canonical ETH assertion: positive deterministic exponential-time rate for deciding every 3-CNF, allowing polynomial input-reading factors.", "scope": "Literature report only. Detailed exposition of externally produced OpenAI results and formalization scope; no new discovery or independent proof certification.", "human_contributions": "Submitter prompted preparation and expansion; no human mathematical authorship or independent proof verification reported.", "verification": "Source checks attributed to supplied report. Integration checks formatting and target identity, not all analytic/algorithmic proofs or Lean dependencies.", "reproducibility": "Supplied manuscript dated2026-10-26, pinned openai/math commit 7ef5c0cc5ca4b0e072b7a050fca09ab24bc6fa3a. Original prompts and browsing logs unavailable; initial/final research-prompt use not documented. Integration expands presentation/link macros and adds metadata without changing mathematical exposition."}
\documentclass[11pt,a4paper]{article}
\usepackage[T1]{fontenc}
\usepackage[utf8]{inputenc}
\usepackage{lmodern}
\usepackage[margin=25mm,headheight=14pt]{geometry}
\usepackage{amsmath,amssymb,amsthm,mathtools}
\usepackage{microtype,booktabs,array,longtable,enumitem}
\usepackage{xurl}
\usepackage[hidelinks]{hyperref}
\usepackage{fancyhdr}
\pagestyle{fancy}
\fancyhf{}
\fancyhead[L]{\small\sffamily WE-SOLVE \; / \; EXTENDED PROOF EXPOSITION}
\fancyhead[R]{\small\sffamily\reportcode}
\fancyfoot[C]{\small\thepage}
\renewcommand{\headrulewidth}{0.3pt}
\setlength{\emergencystretch}{3em}
\setlength{\parskip}{3pt}
\setlist{itemsep=3pt,topsep=5pt}
\allowdisplaybreaks[2]
\numberwithin{equation}{section}
\setcounter{tocdepth}{1}
\newtheorem{theorem}{Theorem}[section]
\newtheorem{proposition}[theorem]{Proposition}
\newtheorem{lemma}[theorem]{Lemma}
\newtheorem{corollary}[theorem]{Corollary}
\theoremstyle{definition}
\newtheorem{definition}[theorem]{Definition}
\theoremstyle{remark}
\newtheorem{remark}[theorem]{Remark}
\newcommand{\R}{\mathbb R}
\newcommand{\C}{\mathbb C}
\newcommand{\Q}{\mathbb Q}
\newcommand{\N}{\mathbb N}
\newcommand{\Z}{\mathbb Z}
\newcommand{\F}{\mathbb F}
\newcommand{\HH}{\mathbb H}
\newcommand{\eps}{\varepsilon}
\newcommand{\dd}{\,\mathrm d}
\newcommand{\poly}{\operatorname{poly}}
\newcommand{\supp}{\operatorname{supp}}
\newcommand{\rank}{\operatorname{rank}}
\newcommand{\tr}{\operatorname{tr}}
\newcommand{\diag}{\operatorname{diag}}
\newcommand{\diver}{\operatorname{div}}
\newcommand{\pc}{p_{\mathrm c}}
\newcommand{\one}{\mathbf 1}
\newcommand{\sourcecommit}{7ef5c0cc5ca4b0e072b7a050fca09ab24bc6fa3a}
\newcommand{\pin}{\texttt{\sourcecommit}}
\newcommand{\repolink}[2]{\href{https://github.com/openai/math/blob/7ef5c0cc5ca4b0e072b7a050fca09ab24bc6fa3a/#1}{#2}}
\newcommand{\mathsource}[2]{\repolink{preprints/Local-smoothing-for-the-Euclidean-wave-equation-October-25-2026/build/sections/#1.tex}{#2}}
\newcommand{\cssource}[2]{\repolink{preprints/The-Exponential-Time-Hypothesis-is-False-October-25-2026/build/#1.tex}{#2}}
\newcommand{\statusbox}[1]{\begin{quote}\small\noindent\textbf{Attribution and verification boundary.} #1\end{quote}}
\hypersetup{pdfauthor={Source-based exposition prepared with GPT-6 Astra Pro},pdfsubject={We-Solve contribution draft; original research attributed to OpenAI}}
\newcommand{\reportcode}{CS-20}
\title{A reported refutation of ETH:\\interpolation, exact counting, and recursive tensor contraction\\[4pt]
\large An extended source-pinned proof exposition for CS-20}
\author{Prepared for \texttt{we-solve} from OpenAI's public research artifacts}
\date{26 October 2026 \quad / \quad Expanded edition}
\begin{document}
\section*{A reported refutation of ETH: interpolation, exact counting, and recursive tensor contraction}
\noindent OpenAI GPT 6 Astra Pro; extended literature report, 26 October 2026.
\section*{Abstract}
This report reconstructs the main proof mechanisms in OpenAI's October 25,
2026 manuscript reporting deterministic subexponential exact counting for every
fixed clause width. It develops the bounded-occurrence interpolation algorithm,
the low-rank replacement construction, balanced incidence schedules, the
Taylor-recovery identity, and the full running-time calculation. It then explains
disjoint sparsification, counting-preserving reductions, the modular parity
lift, deterministic sparse parity tests, exact arithmetic, and polynomial space.
A separate part covers the sharper tensor-contraction route: the Cauchy
deformation, tapered Hermite recovery, valuation-weighted modes, recursive
factorization, and the incidence encoding for arbitrary density. The manuscript's
three quantitative levels are distinguished from the published formal scope.
The final sections derive the precise contradiction to CS-20's ETH formulation,
check the canonical-encoding issue, and delimit the consequences. This is a
detailed literature report about an externally produced proof, not a new
algorithm-discovery claim or an independently completed verification.

\begin{quote}\textbf{Attribution and verification boundary.} OpenAI is credited with the original algorithms, constructions, and
proofs. This report explains the pinned TeX source and formalization documents.
Its displayed source lemmas remain imports from that research. Algebraic
identities and parameter deductions are expanded here, but neither the full
paper nor the entire Lean dependency closure has been independently certified
in this preparation. No Lean execution, practical benchmark, or independent
review is claimed.\end{quote}
\paragraph{Reading route.}
Sections~\ref{sec:encoding}--\ref{sec:simple} give the simplest complete
subexponential route. Sections~\ref{sec:paritylift}--\ref{sec:arithmetic}
explain the improved rational counting bound. Sections~\ref{sec:tensorinterface}--\ref{sec:tensordense}
cover the sharper tensor argument. The concluding sections establish the
exact catalogue implication and its verification boundary.




\section{The exact catalogue claim and the three quantitative levels}
\label{sec:target}
The primary entry is \emph{The Exponential Time Hypothesis for 3-SAT}, CS-20,
with stable handle
\texttt{prb\_a1b2b1f67d0e489da50d97b732d8291d} \cite{catalogue}.
Its formulation asserts a positive deterministic exponential-time rate in the
number $n$ of variables, allowing a polynomial input-reading factor:
\begin{equation}\label{eq:eth}
 \exists c>0:\quad\text{no algorithm decides every $3$-CNF formula in}
 \quad O(2^{cn}\poly(|F|))\text{ time}.
\end{equation}
The source snapshot is the \texttt{openai/math} commit
\begin{center}\texttt{7ef5c0cc5ca4b0e072b7a050fca09ab24bc6fa3a}.\end{center}
The manuscript is dated October 17 and revised October 25, 2026. Source links
below use this commit throughout.

\begin{theorem}[OpenAI's reported main counting theorem]\label{thm:main}
For every fixed $k\ge2$, there is a deterministic algorithm that computes the
number $\#F$ of satisfying assignments of any $k$-CNF with $n\ge20$ declared
variables and $m$ clauses using polynomial space and time
\begin{equation}\label{eq:sharp}
 (n+m)^{O_k(1)}
 2^{O_k(n\log\log n/\sqrt{\log n})}.
\end{equation}
For each fixed $D>0$, inputs with $m\le Dn$ satisfy the sharper bound
\begin{equation}\label{eq:density}
 (n+m)^{O_{k,D}(1)}2^{O_{k,D}(n/\sqrt{\log n})}.
\end{equation}
\end{theorem}
These are the manuscript's main theorem bounds \cite{paper,quant-reduction}.
The documented formal scope instead gives the sufficient but weaker exponent
\begin{equation}\label{eq:rho}
 \rho(n)=\frac{n(\log\log n)^2}{\sqrt{\log n}},\qquad
 (L+1)^{O_k(1)}2^{O_k(\rho(n))},
\end{equation}
where $L$ is the canonical input length \cite{scope}. Small inputs can be
handled separately, or logarithms can be shifted as in the formal target.

There are three useful proof layers, which should not be conflated. The
bounded-occurrence algorithm plus a simple copying reduction already gives
$2^{o(n)}$ time for general fixed-width formulas, although with a weaker
explicit exponent. The sparse-polynomial and parity-test machinery improves
this to \eqref{eq:rho}. The recursive tensor method improves it further to
\eqref{eq:sharp} and \eqref{eq:density}. The first layer already contradicts
ETH. The later layers are quantitative improvements, not missing ingredients
that must be smuggled into the simpler disproof.

Throughout, $k$ is fixed before the input is given. Constants and polynomial
exponents may depend on $k$. Nothing here asserts the same guarantee when the
clause width grows with $n$.

\section{Counting as a polynomial sum, and why this alone is not an algorithm}
\label{sec:encoding}
For each literal on variable $x_v$, write $L_{e,v}(X_v)=X_v$ for a positive
literal and $L_{e,v}(X_v)=1-X_v$ for a negative literal. The satisfaction
polynomial of clause $e$ is
\[
 \Phi_e((X_v)_{v\in I_e})
       =1-\prod_{v\in I_e}(1-L_{e,v}(X_v)).
\]
On Boolean input it is exactly the clause's zero--one indicator. Thus
\begin{equation}\label{eq:count}
 \#F=\sum_{x\in\{0,1\}^n}\prod_e\Phi_e(x|_{I_e}).
\end{equation}
For example,
\[
 F=(x_1\vee x_2\vee x_3)\wedge(\neg x_1\vee\neg x_2\vee x_3)
\]
has factors $1-(1-X_1)(1-X_2)(1-X_3)$ and
$1-X_1X_2(1-X_3)$. Precisely two assignments fail, namely $(0,0,0)$ and
$(1,1,0)$, so the sum is six. This example illustrates the encoding, not a
speedup: evaluating all its summands would still be brute force.

Initially restrict attention to width at most three and at most three
occurrences per variable. Remove tautologies and repeated literals; an empty
clause makes the answer zero. Retain every declared variable. For an unused
variable, add a factor equal to one with that variable as its formal scope.
Then the total number of incidences satisfies
\[
 n\le M:=\sum_e|I_e|\le3n,
\]
and the product in \eqref{eq:count} has degree at most $M$.
The degree bound, rather than a claim that the constraint graph has small
separators, is the starting structural advantage \cite{paper}.

If all clauses lay inside small disjoint variable groups, their counts would
multiply. General instances need not have this property. The source does
something subtler: it replaces clause inputs by polynomially varying linear
forms. At selected interpolation nodes, the modified clauses communicate
through only a few bounded integer measurements. At zero, the original
inputs are restored. Interpolation then recovers the desired count from the
structured evaluations.

\section{Low-degree deformation and the measurement-separation lemma}
\label{sec:design}
For every incidence $(e,v)$ introduce a row of rational polynomials
$P_{e,v}(T)\in\Q[T]^n$ with
\[
 P_{e,v}(0)=\delta_v,\qquad \deg P_{e,v}\le R.
\]
Define
\begin{equation}\label{eq:H}
 H(T)=\sum_{x\in\{0,1\}^n}
       \prod_e\Phi_e\bigl((\langle P_{e,v}(T),x\rangle)_{v\in I_e}\bigr).
\end{equation}
Then $H(0)=\#F$ and $\deg H\le MR\le3nR$. At nonzero values of $T$, the
new clause factors are not necessarily zero or one, or even positive. They
are auxiliary algebraic quantities; correctness is based on interpolation,
not on treating each transformed sum as a satisfiability count.

At a node $t\in\{1,\ldots,Q\}$, a design consists of a partition
$[n]=S_1\sqcup\cdots\sqcup S_r$, with $|S_g|\le2n/r$, an assigned group
$\sigma_e(t)$ for each factor, and integer measurement matrices $C_g$.
The total number of rows is at most $2\tau n$ and every row has absolute
coefficient sum at most $D$. The structural condition is that, for a factor
assigned to $g$, each replacement form can be written
\begin{equation}\label{eq:rowspace}
 \langle P_{e,v}(t),x\rangle
 =\langle P_{e,v}(t)|_{S_g},x_g\rangle+
       \sum_{h\ne g}\langle\lambda_{e,v,h},C_hx_h\rangle.
\end{equation}
All inputs of the same factor must have the same assigned group. Merely
assigning a separate group to each literal would not separate the factor.

Fix $y_g=C_gx_g$ for all groups. Since $x_g$ is Boolean, each coordinate of
$y_g$ lies in $[-D,D]$. Let $A_g(u;y)$ denote the product of the transformed
factors assigned to group $g$, evaluated using local assignment $u$ and the
fixed measurements. Distributivity gives the exact identity
\begin{equation}\label{eq:separate}
 H(t)=\sum_y\prod_{g=1}^r
   \left(\sum_{\substack{u\in\{0,1\}^{S_g}\\C_gu=y_g}}A_g(u;y)\right).
\end{equation}
Every original Boolean assignment occurs once: it has one measurement tuple
and one choice of local assignment in every group. Empty measurement fibers
contribute zero. This establishes correctness without assuming independence
of the original clauses.

Enumerating the tuples and local assignments takes, up to polynomial factors,
\begin{equation}\label{eq:sep-cost}
 (2D+1)^{2\tau n}2^{2n/r}.
\end{equation}
The procedure can keep one tuple and one local sum at a time. It need not
store a table indexed by all measurement tuples.

Both the number of measurements and their coefficient size matter. A single
row $(1,2,4,\ldots,2^{n-1})$ has rank one but records the entire Boolean
assignment. Thus low rank alone would not imply a small enumeration domain.
The explicit row-$\ell^1$ bound is an indispensable part of the design.

\section{Recovering Taylor coefficients without destroying separation}
\label{sec:taylor}
The $Q$ values $H(t)$ do not determine a polynomial whose degree can be
$3nR$. The source uses Hermite interpolation, with
\[
 q=\left\lfloor\frac{3nR}{Q}\right\rfloor,\qquad J=q+1.
\]
Since $QJ>3nR$, the first $J$ Taylor coefficients at each node suffice.
Indeed, two polynomials with the same data have a difference divisible by
$\prod_{t=1}^Q(T-t)^J$, of degree exceeding the allowed degree.
The issue is how to compute these coefficients using designs that hold only
at the nodes themselves.

Fix $t$, put $\Delta P(z)=P(t+z)-P(t)$, and introduce a second formal vector
$X$, distinct from the Boolean summation variable $x$. Set
\[
 G_x(X,z)=\prod_e\Phi_e\Bigl(
 \bigl(\langle P_{e,v}(t),x\rangle+
 \langle\Delta P_{e,v}(z),X\rangle\bigr)_{v\in I_e}\Bigr).
\]
Setting $X=x$ recovers the shifted product, so
$H(t+z)=\sum_xG_x(x,z)$. Each occurrence of $X$ in this expression carries
at least one factor of $z$. Consequently
\begin{equation}\label{eq:taylor-degree}
 \deg_X[z^j]G_x(X,z)\le j.
\end{equation}

Work modulo $z^J$ and multilinearize in $X$, replacing positive powers of
$X_i$ by $X_i$. This does not change values on the Boolean cube. The result
can be written
\[
 G_x(X,z)\equiv\sum_{|E|\le q}a_E(x,z)\prod_{i\in E}X_i\pmod{z^J}.
\]
Evaluating at the Boolean vector $\one_U$ sums the coefficients whose supports
are contained in $U$. Subset inversion gives
\[
 a_E(x,z)=\sum_{U\subseteq E}(-1)^{|E|-|U|}G_x(\one_U,z)
                 \pmod{z^J}.
\]
On the diagonal $X=x$, the monomial for $E$ is one exactly when all bits in
$E$ are one. Therefore
\begin{equation}\label{eq:jetidentity}
 H(t+z)\equiv
 \sum_{|E|\le q}\sum_{U\subseteq E}(-1)^{|E|-|U|}
 \sum_{\substack{x\in\{0,1\}^n\\x_i=1\ (i\in E)}}
                         G_x(\one_U,z)\pmod{z^J}.
\end{equation}
This is a full algebraic identity, not a heuristic approximation to the
Taylor series.

For fixed $U$, the correction $\langle\Delta P(z),\one_U\rangle$ no longer
depends on $x$. Thus the original design still separates the dependence on
$x$. The restrictions $x_i=1$ are local to the groups. Formula
\eqref{eq:separate} applies over the truncated ring $\Q[z]/(z^J)$, because
its proof uses only distributivity. The algorithm evaluates the nested sums
without ever constructing a multivariate expansion in all $X_i$.

The number of pairs $(E,U)$ is
\begin{equation}\label{eq:taylorcount}
 C_q=\sum_{j=0}^{\min(q,n)}2^j\binom nj.
\end{equation}
For $1\le q\le n/2$,
$C_q\le(q+1)(2en/q)^q$, so its logarithm is
$O(n(R/Q)\log(2Q/R))$. Multiplying by \eqref{eq:sep-cost} shows that the
exponent is controlled by
\begin{equation}\label{eq:threecosts}
 n\left[\frac RQ\log\frac{2Q}{R}
             +\tau\log(2D+1)+\frac1r\right].
\end{equation}
The construction must make \emph{all three} terms small simultaneously.

\section{The auxiliary matrix: spectral concentration and sparse mismatches}
\label{sec:matrix}
The source constructs the replacement rows in an auxiliary coordinate space
\[
 \mathcal B=([r]\times[r])^Q,
 \qquad D=|\mathcal B|=r^{2Q}.
\]
Write a coordinate as $u=((a_u(s),d_u(s)))_{s=1}^Q$. At node $t$, coordinate
$w$ belongs to group $a_w(t)$, while row $u$ is assigned to group $d_u(t)$.
The mismatch set is
\[
 E_{u,w}=\{t:a_w(t)\ne d_u(t)\}.
\]
The aim is to find a matrix $M_0$ that is the identity on a large subspace,
but whose nonzero entries have small mismatch sets.

\subsection{Commuting swaps make the two requirements compatible}
Let $\Pi_s$ swap the two symbols in coordinate $s$. These matrices commute
and satisfy $\Pi_s^2=I$. Define
\[
 V_0=\prod_{s=1}^Q\Pi_s,\qquad
 W_0=\frac12\sum_{s=1}^Q(I-\Pi_s),\qquad M_0=V_0f(W_0).
\]
If $f$ has degree at most $R-1$, expanding it gives products of at most
$R-1$ swaps. Multiplication by $V_0$ means that all coordinates are swapped
except perhaps those $R-1$ positions. A complete swap sends row $u$ to a
column with $a_w(t)=d_u(t)$ everywhere. Thus
\begin{equation}\label{eq:supportmatrix}
 M_0[u,w]\ne0\quad\Longrightarrow\quad |E_{u,w}|\le R-1.
\end{equation}

For one pair of symbols, the swap has an orthogonal basis consisting of
$e_{aa}$ and $e_{ab}+e_{ba}$ with eigenvalue $+1$, and $e_{ab}-e_{ba}$ with
eigenvalue $-1$. Tensoring these bases across the $Q$ positions diagonalizes
all swaps. If a basis vector has $j$ negative factors, its eigenvalues under
$W_0,V_0,M_0$ are respectively
\[
 j,\qquad(-1)^j,\qquad(-1)^jf(j).
\]
Under uniform choice of a tensor basis vector, $j$ has distribution
$\operatorname{Bin}(Q,p)$, with $p=(r-1)/(2r)$.

Choose an interval $\mathcal I$ of $R$ consecutive integers near $Qp$ and
interpolate $f(j)=(-1)^j$ for $j\in\mathcal I$. This gives degree at most
$R-1$ and makes $M_0$ fix every central basis vector. A concrete formula when
$\mathcal I=\{A_0,\ldots,A_0+R-1\}$ is
\[
 f(z)=(-1)^{A_0}\sum_{j=0}^{R-1}(-2)^j\binom{z-A_0}{j}.
\]
The degree stays small because it is the width of the central interval, not
the total number of nodes, that matters.

Let $C$ have as its rows the exceptional tensor basis vectors. Its entries
are $0,1,-1$, and its kernel is spanned by the central basis. A binomial tail
bound gives
\begin{equation}\label{eq:exceptional}
 \#\operatorname{rows}(C)
 \le2D e^{-R^2/(8Q)}.
\end{equation}
Since $I-M_0$ vanishes on $\ker C$, every row of $I-M_0$ belongs to the row
space of $C$. Its output is therefore determined by the few measurements
$C\zeta$.

\subsection{A polynomial kills only the mismatched dependencies}
For each nonzero $M_0[u,w]$, set
\[
 \varphi_{u,w}(T)=\prod_{s\in E_{u,w}}(1-T/s).
\]
Then define the replacement form by
\begin{equation}\label{eq:auxrow}
 \langle P_u(T),\zeta\rangle
 =((I-M_0)\zeta)_u+
   \sum_wM_0[u,w]\varphi_{u,w}(T)\zeta_w.
\end{equation}
At zero, all multipliers equal one and the expression is $\zeta_u$.
By \eqref{eq:supportmatrix}, its degree is at most $R-1$.
At $T=t$, every term with $a_w(t)\ne d_u(t)$ vanishes. The surviving part
uses only the assigned group, while the first term is determined by $C\zeta$.
This proves the desired auxiliary measurement property.

The construction explains the unusual conjunction of requirements. Spectral
concentration makes the correction low rank; low degree in the swap matrices
bounds the mismatches; and a product of roots at those mismatches turns the
matrix into a low-degree interpolation family. None of these three facts
alone would give the algorithm \cite{paper}.

\section{Placing arbitrary incidences into the auxiliary coordinates}
\label{sec:balance}
The auxiliary rows must now accommodate the actual formula, including repeated
uses of the same variable. Give each variable a schedule $a_v\in[r]^Q$ and
each factor a schedule $d_e\in[r]^Q$. Its incidence $(e,v)$ receives label
$u(e,v)=(a_v,d_e)\in\mathcal B$.

For the more general construction, assume factor scopes have size at most $w$
and each variable occurs in at most $b$ scopes. Let $M$ be total incidences,
with $n\le M\le bn$. The source chooses schedules deterministically so that
\begin{equation}\label{eq:balance}
 |S_{t,g}|\le2n/r,\qquad
 \max_u\#\{(e,v):u(e,v)=u\}\le2M/D,
\end{equation}
provided $2(b+w)D^2\le M$ and $2Qr^2\le n$.

Here is the conditional-expectation argument. Under independent uniform
schedules, let $c_u$ be label loads and $n_{t,g}$ group sizes. Expanding squared
deviations gives
\[
 \mathbb E\sum_u(c_u-M/D)^2\le M(b+w),\qquad
 \mathbb E\sum_{t,g}(n_{t,g}-n/r)^2\le Qn.
\]
Pairs of incidences with disjoint variables and factors have zero covariance;
only shared incidences contribute to the first bound. Fix schedules one by
one while not increasing the normalized sum of these two potentials. Each
conditional expectation concerns at most four schedules per pair of
indicators and is computable by explicit finite enumeration in polynomial
size. The final potentials are at most twice their respective scales, giving
\[
 c_u\le M/D+\sqrt{2M(b+w)},\qquad
 n_{t,g}\le n/r+\sqrt{2Qn}.
\]
The displayed numerical assumptions imply \eqref{eq:balance}.

Use $h_0=\max_uc_u$ copies of $\mathcal B$ to place incidences with the same
label into different copies. Put the original variable $X_v$ at its incidence
position, and zero in unused positions. Repeated appearances of $X_v$ still
have the same value: these copies are not additional independent Boolean
variables. Since $h_0D\le2M$, the total coordinate budget remains polynomial.

Apply \eqref{eq:auxrow} in the appropriate copy. Measurements can involve
several actual variable groups, so split each measurement into its restrictions
to those groups. This step is essential: conditioning only on a sum of two
groups could continue to couple them. If the exceptional fraction in
\eqref{eq:exceptional} is at most $\tau/(rb)$, the total number of split rows
is at most
\[
 rh_0D\frac{\tau}{rb}\le2\tau n.
\]
Each row still has integer coefficients with absolute sum at most $D$.
Combining repeated occurrences cannot increase that bound. Thus the transferred
rows satisfy the design of Section~\ref{sec:design}.

\section{Parameter choices and the bounded-occurrence running time}
\label{sec:parameters}
Put $\lambda=\lceil\log_2n\rceil$ and
$\nu=\lceil\log_2\lambda\rceil$. For sufficiently large $n$, the source uses
parameters of the following orders:
\begin{align*}
 r&=\left\lfloor\frac{\sqrt\lambda}{\nu^2}\right\rfloor,
 &Q&=\left\lfloor\frac{\lambda}{128\nu}\right\rfloor,\\
 R&=(\beta+3)\lceil\sqrt\lambda\rceil,
 &\tau&=\frac1{r\lambda},\qquad D=r^{2Q}.
\end{align*}
Here $\beta$ is a fixed bound such that $b,w\le\lambda^\beta$; the
width-three, occurrence-three case needs only a fixed choice.
One obtains $\log_2D\le\lambda/128$, so $D\le(2n)^{1/128}$.
This ensures the balancing inequalities because their left sides have a
strictly smaller power of $n$ than their right sides.

The ratio $R/Q$ is $O(\nu/\sqrt\lambda)$ and tends to zero. The interval of
$R$ central binomial levels therefore fits inside $\{0,\ldots,Q\}$ for large
$n$. Its tail is exponentially small in $R^2/Q$, which is a sufficiently large
fixed multiple of $\nu$; this is enough to make the exceptional fraction at
most $\tau/(rb)$. All these requirements can be checked directly from the
integer parameters. The finitely many sizes failing the requirements can be
handled by exhaustive enumeration.

Substitution into \eqref{eq:threecosts} gives
\[
 \frac RQ\log\frac{2Q}{R}=O\left(\frac{\nu^2}{\sqrt\lambda}\right),
 \qquad \tau\log(2D+1)=O(1/r),
 \qquad \frac1r=O\left(\frac{\nu^2}{\sqrt\lambda}\right).
\]
Hence bounded-occurrence counting takes
\begin{equation}\label{eq:bounded}
 n^{O(1)}2^{O(n(\log\log n)^2/\sqrt{\log n})}
\end{equation}
bit operations, once the arithmetic bounds below are included.
The matrices, schedules, polynomial rows, and Hermite systems are only
polynomially large. The subexponential factor comes from the three explicit
enumerations, not from an uncounted construction of enormous matrices.

\section{Disjoint sparsification: reducing density without double counting}
\label{sec:sparsification}
To pass from a special class to all fixed-width formulas, a decision-only
sparsification statement is insufficient for exact counting: overlapping
solution sets would be counted repeatedly. The source proves a disjoint
branching version with explicit dependence on its sparsity parameter
\cite{paper}.

Normalize the formula into an antichain of literal sets by deleting tautologies,
duplicates, and any clause containing a smaller clause. Retain the declared
variable set rather than silently deleting unused variables. Fix an integer
$K\ge2$ and thresholds $\theta_j=K^j$. If a nonempty proper literal set $H$
is contained in too many clauses of a width $c$, branch on whether $H$ is
satisfied. Choose the smallest eligible width and, within it, a core $H$ of
largest size, with deterministic tie-breaking.

For selected clauses $H\cup P_1,\ldots,H\cup P_t$, and conjunction $G$ of
other clauses, the two children are
\begin{align}
 F_A&=G\wedge H,\label{eq:branchA}\\
 F_B&=G\wedge\bigwedge_{j=1}^tP_j
                   \wedge\bigwedge_{\ell\in H}\{\neg\ell\}.
                                                        \label{eq:branchB}
\end{align}
In the first child at least one literal of $H$ is true, so all selected clauses
are automatically satisfied. In the second, every literal of $H$ is false,
so each selected clause reduces to its petal $P_j$. These cases are exhaustive
and disjoint. Petals need not be pairwise disjoint: this is a branching identity,
not an unsupported sunflower assumption.

The antichain condition and choice of core ensure that newly introduced
smaller clauses genuinely survive normalization. Mark such clauses to track
the recursion. The source bounds the number of marked width-$j$ clauses
containing a fixed smaller literal set by a multiple of $K^{j-i}$.
If $T_j$ counts relevant insertions, the resulting recurrence has the form
\[
 T_1\le2n,\qquad
 T_j\le4nK^{j-1}+2\sum_{i<j}K^{j-i}T_i
       \le4(3K)^{j-1}n.
\]
Thus the branch depth is at most $C_kn$, where $C_k$ is a fixed polynomial
in $K$, and the number of more expensive B-choices on a path is at most
$\Gamma_kn/K$.

Encoding a path by its length and the positions of its B-choices bounds the
tree by a polynomial factor times
\[
 2^{(\Gamma_kn/K)\log_2(4C_kK)}.
\]
Choose $K=O_k(\sigma^{-1}\log(2/\sigma))$ so this exponent is at most
$\sigma n$. Every leaf then has variable-occurrence bound
\begin{equation}\label{eq:sparseB}
 B=O_k\left((\sigma^{-1}\log(2/\sigma))^{k-1}\right)
\end{equation}
and at most $Bn$ nonempty clauses. The tree can be traversed depth first in
polynomial space, and its solution sets partition those of the original
formula. The explicit parameter dependence permits $\sigma$ to vary slowly
with $n$; it is not merely a theorem with an opaque constant for each fixed
$\sigma$.

\section{The simpler complete route to a subexponential algorithm}
\label{sec:simple}
Before introducing the sharper machinery, it is useful to close a complete
logical route from \eqref{eq:bounded} to an ETH refutation.
For a sparse fixed-width leaf $G$, compute the disjunction in each long clause
using binary OR gates. A gate $z=u\vee v$ is enforced by
\[
 (\neg u\vee z)\wedge(\neg v\vee z)\wedge(u\vee v\vee\neg z).
\]
Every assignment to the gate inputs determines exactly one valid output.
Requiring the final gate output to be one therefore gives a width-three
encoding with a unique extension of every satisfying original assignment.
This is count preserving, not merely satisfiability preserving.

If a variable has $s\ge2$ occurrences, replace them by copies
$x_1,\ldots,x_s$ and add the cycle
\[
 (\neg x_1\vee x_2)\wedge\cdots\wedge
 (\neg x_s\vee x_1).
\]
A cyclic sequence of Boolean implications forces all copies to agree. Each
copy appears in one inherited clause and two cycle clauses, so at most three
times. Variables with zero or one occurrence remain, preserving their counting
contribution. The resulting instance has $N=O_k(Bn)$ variables.

Applying \eqref{eq:bounded} to every leaf gives exponent
\[
 O_k\left(\sigma n+
       Bn\frac{(\log\log n)^2}{\sqrt{\log n}}\right)
\]
when $B$ is polylogarithmic. For example, taking
$\sigma\asymp(\log n)^{-1/(2k)}$ and using \eqref{eq:sparseB} gives
\begin{equation}\label{eq:simple-general}
 (n+m)^{O_k(1)}
 2^{O_k(n(\log\log n)^{k+1}/(\log n)^{1/(2k)})}
 =(n+m)^{O_k(1)}2^{o(n)}.
\end{equation}
This already solves arbitrary fixed-width SAT in subexponential time, by
checking whether the exact count is positive. No parity-test lemma or tensor
contraction is needed for this qualitative conclusion.

It would, however, be incorrect to claim the stronger exponent $O(\rho(n))$
from this copying route alone. The factor $B$ multiplies the number of variables.
The next construction avoids that blowup by staying on the original Boolean
variables.

\section{A modular parity lift with a global degree bound}
\label{sec:paritylift}
Consider $a=O(n)$ equations over $\F_2$, each represented by a sparse polynomial
of fixed degree at most $k$. Lift equation $i$ to the ordinary integer sum
$L_i(X)$ of its listed monomials. For Boolean $x$, the equation holds exactly
when $L_i(x)$ is even.

Set $\ell=n+1$. The source uses the integer polynomial
\begin{equation}\label{eq:paritypoly}
 p_\ell(Y)=(1-Y)^\ell
        \sum_{j=0}^{\ell-1}\binom{\ell+j-1}{j}Y^j.
\end{equation}
The truncated binomial series gives $p_\ell(Y)\equiv1\pmod{Y^\ell}$, while
the prefactor gives $p_\ell(Y)\equiv0\pmod{(1-Y)^\ell}$. Therefore
\[
 p_\ell(y)\equiv
 \begin{cases}1&y\text{ even},\\0&y\text{ odd},\end{cases}
 \pmod{2^\ell}.
\]
The precision is large enough because a count is between zero and $2^n$, so
its least nonnegative residue modulo $2^{n+1}$ determines it exactly.

A direct product of these parity polynomials would have degree potentially
quadratic in $n$. The key is to truncate the product globally. Put
$h(Y)=Y(Y-1)$, which is even on every integer. Repeated division by this monic
quadratic gives
\begin{equation}\label{eq:paritybasis}
 p_\ell(Y)=\sum_{j=0}^{\ell-1}e_j(Y)h(Y)^j,
 \qquad e_j\in\Z[Y],\quad\deg e_j\le1.
\end{equation}
Introduce one common marker $V$ and work in $\Q[V]/(V^\ell)$ with factors
\[
 B_i(X,V)=\sum_{j=0}^{\ell-1}e_j(L_i(X))h(L_i(X))^jV^j.
\]
After multiplying, only terms with total marker degree below $\ell$ remain.
Every discarded term has at least $\ell$ copies of an even integer when
$X=x$ is Boolean, so vanishes modulo $2^\ell$. The retained incidence-degree
bound is
\begin{equation}\label{eq:totaldegree}
 k\bigl(a+2(\ell-1)\bigr)=O(n).
\end{equation}
The word ``total'' is essential: truncating each factor separately would not
impose this linear bound on their product.

Let $\mathcal S_\ell$ sum the retained coefficients of $V$. Then the number of
common zeros is the least residue of
\[
 \mathcal S_\ell\left(\sum_{x\in\{0,1\}^n}\prod_iB_i(x,V)\right)
 \pmod{2^\ell}.
\]
The coefficient-sum map is applied \emph{after} multiplication in the
truncated ring. It is not a ring homomorphism, and it cannot be treated as
substitution $V=1$ in the quotient: the relation $V^\ell=0$ would then be
violated. This distinction is important for exactness \cite{paper}.

At deformed rational inputs the modular parity interpretation need not hold.
The interpolation algorithm therefore computes over the rationals, recovers
the value at $T=0$, and only then takes the integer residue. It does not divide
by arbitrary even denominators modulo a power of two.

\section{Counting sparse polynomial systems on the original variables}
\label{sec:sparsecounter}
Suppose the number of equations is at most $cn$ and each has at most
$(\log n)^c$ monomials of degree at most $k$, for fixed $k,c$.
Its scope is therefore polylogarithmic, although a few variables might occur
in many scopes. Let $M$ be total incidences, adding neutral factors for unused
variables as needed. Choose a threshold comparable to
$(M/n)\sqrt{\log n}$. Fewer than $n/\sqrt{\log n}$ variables exceed it.
Enumerate their values, paying only $2^{O(n/\sqrt{\log n})}$ branches.

On each branch, both scope size and occurrence count are bounded by fixed
powers of $\log n$. The design construction of Sections~\ref{sec:matrix}--
\ref{sec:parameters} therefore applies with a suitable fixed $\beta$.
Use one replacement form per variable in a test and reuse it throughout that
test. The bound \eqref{eq:totaldegree} implies that the deformed polynomial
has degree $O(nR)$ in $T$, despite the polylogarithmic number of variables
read by a factor.

The same Taylor identity applies over the coefficient algebra
$\Q[V,z]/(V^\ell,z^J)$. It tracks only the two bounded degrees, not all
monomials in the original variables. A factor is evaluated directly from its
listed sparse polynomial; no truth table over its full scope is required.
The branch count is recovered by rational Hermite interpolation, followed by
$\mathcal S_\ell$ and reduction modulo $2^\ell$. The original precision
$\ell=n+1$ remains valid after branching.

The source thereby obtains a deterministic counter for \emph{arbitrary}
supplied sparse polynomial systems in
\begin{equation}\label{eq:sparsecounter}
 2^{O_{k,c}(\rho(n))}\text{ time and }n^{O_{k,c}(1)}\text{ space}.
\end{equation}
This generality is what makes the next derandomization noncircular. The
counter is already valid for systems which accept unwanted assignments; it
does not require a promise that the system represents exactly the solutions
of a CNF formula.

\section{Sparse parity tests: first exclude low-weight violation patterns}
\label{sec:expander}
For a sparse CNF with $m\le Bn$ clauses, let
$v_e(x)$ be one when clause $e$ is violated and zero otherwise. A parity test
is a linear equation in this violation vector over $\F_2$.
Every satisfying assignment passes all such tests. The task is to choose only
$O(n)$ short tests so that no nonsatisfying assignment passes them all.

A random short test is unlikely to see a violation vector supported on just
one clause. The source first constructs a parity-check system whose nonzero
kernel vectors have large support. Write
$\mu_0=\lfloor\log_2m\rfloor+1$ and use, for example, the source's scales
\[
 P=20000B^2\mu_0^2,\qquad D_0=320000B^2\mu_0^3.
\]
It constructs at most $3n$ initial rows, each supported on at most $D_0$
clauses, such that every surviving nonzero violation vector has weight at
least $m/P$ \cite{paper}.

\subsection{The graph and the local code}
The construction is an edge-based code on a bipartite expander multigraph.
Both sides are copies of a binary vector group. An edge at port $j$ joins
$u$ on the left to $u\oplus s_j$ on the right. Parallel edges are distinct
ports. Choose shifts $s_j$ so that every nonconstant binary character has
small average on them. The character basis diagonalizes the adjacency
operator, giving the mixing bound
\begin{equation}\label{eq:mixing}
 e(U,V)\le \frac{d_G|U||V|}{q_G}
              +\eps d_G\sqrt{|U||V|}.
\end{equation}
The number of vertices $q_G$ and degree $d_G$ are chosen so the graph has
between $m$ and a constant multiple of $m$ edges.

At every vertex impose a small binary linear code with relative distance
$\delta$, where $\delta=2\eps$ and both are inverse-polylogarithmic in the
application. A nonzero global codeword determines sets $U,V$ of vertices
incident to its nonzero edges. Local distance gives at least $\delta d_G$
nonzero edges at each active vertex, while \eqref{eq:mixing} bounds how many
edges can lie between $U$ and $V$. Combining the two implies
\[
 \sqrt{|U||V|}\ge(\delta-\eps)q_G,
 \qquad
 \text{weight}\ge\delta(\delta-\eps)q_Gd_G.
\]
This is the asserted lower support threshold after the parameter choices.
Shorten the code to exactly $m$ coordinates by fixing extra edge bits to zero;
shortening preserves the minimum-distance lower bound on nonzero words.

\subsection{Why this construction is deterministic and affordable}
The source establishes existence of the shifts by a concentration estimate
for every nontrivial character, and existence of the local code by a union
bound over its low-weight words. It then searches deterministically for
objects satisfying the finite stated tests. The searched descriptions have
size polynomial in $B$ and $\log m$, so the work is
$2^{\poly(B,\log m)}$ and can be performed sequentially in polynomial space
in the main input size.

When $B$ is a fixed power of $\log n$, this is
$2^{(\log n)^{O_k(1)}}$, negligible beside $2^{O_k(\rho(n))}$ for sufficiently
large $n$. It is not necessarily polynomial-time construction, and describing
it that way would undercount the work. Its subexponential construction cost
is explicitly absorbed in the final bound.

\section{Greedy deterministic tests: eliminating every remaining false positive}
\label{sec:greedy}
After the initial code checks, any nonsatisfying assignment still passing has
at least $m/P$ violated clauses. Consider a candidate test obtained from $P$
independent clause indices, each with an independent coefficient bit. Repeated
indices are allowed, with repeated selected occurrences cancelling modulo two.
The chance of hitting the violation support is at least $1-e^{-1}>1/2$.
Conditional on a hit, one coefficient bit makes the parity odd with probability
$1/2$, irrespective of the other coefficient bits. Thus the rejection
probability is at least $1/4$.

For a current test matrix $A$, write
\[
 N(A)=\#F+D(A),
\]
where $D(A)$ is the number of nonsatisfying assignments still accepted.
Averaging over candidates gives a candidate $h$ with
$D(A;h)\le(3/4)D(A)$. The algorithm does not know $\#F$ or $D(A)$.
Instead, it computes $N(A;h)$ for every candidate using the already established
sparse-polynomial counter and selects one with minimum count.
Because every genuine solution passes every candidate, $\#F$ is the same
additive baseline for all comparisons; minimizing $N$ also minimizes $D$.

After $3n$ such choices,
\[
 D\le2^n(3/4)^{3n}=(27/32)^n<1.
\]
Since $D$ is a nonnegative integer, it is zero. There are at most $6n$ rows
in total, combining the initial and added tests. Each test becomes a degree
at most $k$ polynomial equation with at most
$2^k\max\{P,D_0\}$ listed monomials, because each clause violation polynomial
has at most $2^k$ monomials.

The candidate family has at most $(2m)^P$ members per round. Consequently the
number of counter calls is at most $3n(2m)^P+1$, again
$2^{(\log n)^{O_k(1)}}$ for polylogarithmic $B$. Calls are processed one at a
time. This proves a deterministic count-preserving reduction on the original
$n$ variables. There is no circular appeal to a SAT oracle: each comparison
is a concrete call to \eqref{eq:sparsecounter}, which is valid even for the
rejected candidates' larger solution sets.

\section{Assembling the rational route and accounting for bits and space}
\label{sec:arithmetic}
Choose the disjoint sparsification parameter on the order of
$(\log\log n)^2/\sqrt{\log n}$. Then $B$ is polylogarithmic, the number of
leaves is $2^{O(\rho(n))}$, and the parity-query cost at each leaf is also
$2^{O_k(\rho(n))}$. Multiplication of these counts adds their exponents, so
the overall bound is \eqref{eq:rho}, up to a polynomial input-size factor.
This conclusion still needs exact arithmetic and space bounds; arithmetic
operations on arbitrarily large rational numbers cannot be charged unit cost.

\subsection{Controlled denominators, not an exponential denominator product}
The interpolation rows admit a common denominator built from factorials,
for example $(R-1)!Q!$ before the fixed-degree substitutions. Its bit length
is polynomial. Coefficient growth in the swap-polynomial construction is
bounded by $2^{O(Q\log Q)}$, also polynomial in bit length.

For the sparse-polynomial lift, choose a common denominator for every
substituted factor, then a common denominator for the entire product.
The crucial implementation invariant is that the \emph{same} denominator is
used for all summands at a given stage. Exponentially many additions then
increase numerator bit length by the logarithm of their number; they do not
multiply the denominator afresh at every addition.

The rings $\Q[V]/(V^\ell)$ and
$\Q[V,z]/(V^\ell,z^J)$ have polynomially many stored coordinates.
The coefficients $e_j$ in \eqref{eq:paritybasis} have polynomial bit size,
as follows by repeated division by a fixed monic quadratic. Products and
truncations therefore have polynomial coordinate bit bounds under the
explicit common-denominator invariant.

\subsection{Hermite reconstruction is an exact polynomial-size linear system}
If the coefficients of $H$ are unknowns, the equation for Taylor coefficient
$j$ at node $t$ has coefficient $\binom uj t^{u-j}$ in column $u\ge j$.
The system is an integer confluent Vandermonde system of polynomial dimension,
with entries of polynomial bit length. Its nonsingularity is the divisibility
argument of Section~\ref{sec:taylor}.

After clearing the common denominator of the data, fraction-free elimination
or a determinant-ratio computation gives the desired constant coefficient.
Bounds on minors give polynomial bit lengths; one does not expand a determinant
by enumerating all permutations. At $T=0$ the recovered quantity is integral.
Only then is the marker-coefficient sum reduced modulo $2^{n+1}$ to obtain
the exact count.

\subsection{Streaming every large enumeration}
The algorithm processes sparsification branches, heavy-variable assignments,
parity candidates, interpolation nodes, subset pairs $(E,U)$, measurement
tuples, and local Boolean assignments sequentially. The data retained at any
time are polynomially many counters, matrices, polynomial coefficients, and
running sums. The disjoint leaf counts add to at most $2^n$, so the final
accumulator needs only $n+1$ bits.

Polynomial space is thus a concrete storage assertion, not a consequence of
subexponential time. The source's formal statement strengthens this further
by bounding all visited tape cells at every prefix of a primitive-machine
execution \cite{scope,challenge}.

\section{The sharper route: Boolean tensor contraction and fixed density}
\label{sec:tensorinterface}
The sharper manuscript bounds use a second algorithmic engine, not simply a
better estimate of \eqref{eq:taylorcount}. A Boolean tensor network is a
bipartite multigraph with one bit on each edge. A vertex tensor assigns a
coefficient-ring element to every assignment of its incident port bits. Its
contraction is
\[
 Z=\sum_{x\in\{0,1\}^{E(G)}}\prod_{v\in V(G)}T_v(x|_{\delta(v)}).
\]
Parallel edges have distinct ports. Size is measured by the number of edges,
not by a potentially exponential explicit truth table for every tensor.

The source proves a contraction theorem for $O(s)$ edges, maximum degree
at most $\lceil s^{9/10}\rceil$, and entries in
\[
 \mathcal B_{P_0,\ell}=(\Z/2^{P_0}\Z)[V]/(V^\ell),
 \qquad P_0,\ell=O(s).
\]
There are polynomially many vertices. Each entry is specified by a procedure
using polynomial space and at most a polynomial factor times
$2^{O(s^{9/10})}$ time. Under these hypotheses the contraction is computed
in polynomial space and
\begin{equation}\label{eq:tensorbound}
 (s+\mathcal L)^{O(1)}2^{O(s/\sqrt{\log s})}
\end{equation}
bit operations, where $\mathcal L$ is the description length
\cite{quant-reduction,quant-contraction}.
The entry-access assumption is part of the theorem; arbitrary tensors given
only by unspecified oracles are not covered.

For a fixed-width CNF, make a clause tensor for each clause and equality
tensors to identify the ports representing the same variable. High-degree
equality constraints can be replaced by trees of trivalent equality tensors.
Every consistent setting of the external ports extends uniquely through such
a tree. Subdivide every edge once by a binary equality tensor to obtain a
bipartite network. The resulting graph has $O_k(n+m)$ edges and bounded
maximum degree.

Apply the contraction theorem with $s=n+m$, $P_0=n+1$, and $\ell=1$.
Unused declared variables contribute a factor $2$ each; multiply by their
combined factor before taking the least residue. The count is at most $2^n$,
so reduction modulo $2^{n+1}$ is exact. If $m\le Dn$, then $s=O_D(n)$,
and \eqref{eq:tensorbound} gives \eqref{eq:density}. This reduction has no
variable blowup depending on a polylogarithmic sparsity parameter.

The remaining sections explain why the contraction theorem has this strong
bound, and how sparse parity systems can be encoded with only
$O(n\log\log n)$ wires for the arbitrary-density result.

\section{Exact coefficient arithmetic and balanced packets}
\label{sec:tensorpackets}
\subsection{A finite extension supplies enough interpolation points}
Let $\lambda=\lceil\log_2s\rceil$, $e=4\lambda$, and choose a monic integer
polynomial $f$ of degree $e$ whose reduction modulo two is irreducible. The
source finds one by finite deterministic enumeration, polynomial in $s$.
Work in
\[
 \mathcal O=\Z_{(2)}[U]/(f),\quad
 \mathcal K=\Q[U]/(f),\quad
 \mathcal A=\mathcal O[V]/(V^\ell),\quad
 \mathcal A_p=\mathcal A/(2^p).
\]
Here $\Z_{(2)}$ allows odd denominators. Nonzero residues modulo two are units,
because the residue ring is the field $\F_{2^e}$. The valuation of a nonzero
scalar in $\mathcal K$ is the minimum two-adic valuation of its coordinates
in the basis $1,U,\ldots,U^{e-1}$; reduction to the residue field shows it is
multiplicative.

Binary-coordinate lifts of distinct residue-field elements have unit
differences. Thus there are sufficiently many points for ordinary Lagrange
interpolation modulo $2^p$ without dividing by nonunits. Arithmetic stores
$e\ell$ coordinates, each with $p$ bits. Unit inverses are computed by exact
linear algebra or lifting from characteristic two. These operations have
polynomial bit complexity at the precisions used by the algorithm.

When precision increases recursively, each tensor uses its original
creation-time coordinate representatives. This convention is important:
canonical lifting of residues is not a ring homomorphism. The proof does
not treat it as one. Instead it fixes entrywise lifts once for a call's
deformation, and all its evaluations refer to those same lifted tensors.

\subsection{Recursive scales and collision control}
The root size parameter $s$ remains fixed. A call with $M$ wires is a leaf
when $M\le H$, where $H\asymp s/\sqrt\lambda$. At active depth $d$, use a
fixed number of groups $r=64$ and parameters
\[
 u_d=\left\lfloor\frac{\eps_0\lambda}{2^d}\right\rfloor,
 \qquad a_d=2\sum_{j=0}^du_j,
 \qquad L=2^{a_d},\qquad D_{\rm arr}=r^{2a_d},
\]
with a sufficiently small absolute constant $\eps_0$.
The source arranges $M\le Cs8^{-d}$ and
$L,D_{\rm arr}\le s^{1/1000}$. At active calls, $u_d$ is still a large power
of $\lambda$, while the maximum vertex degree remains at most
$\Delta=\lceil s^{9/10}\rceil$.

Assign vertices on each side labels in $[r]^{a_d}$. The second-moment
calculation for pair-label loads has expected squared deviation at most
$2M\Delta$: only pairs of edges sharing a vertex can contribute covariance.
Conditional expectation yields a deterministic assignment with the same
bound. Padding each pair-label class to a common multiple of $L$ then gives
at most $3M/L$ packets of $L$ positions.

Each packet fixes the unordered pair of group symbols in every coordinate.
When the pair has distinct symbols, the corresponding bit of the position
index specifies its orientation. In any coordinate, each group has at most
$3M/r$ real ports on either side. The packing is possible because
$M\ge2D_{\rm arr}^2\Delta$ and $M\ge LD_{\rm arr}$ at active scales.
The unused padded positions carry no real ports. This combinatorial step
supplies both the bit structure used below and the eventual factor-eight
shrinkage of recursive networks \cite{quant-contraction}.

\section{The Cauchy deformation and tapered Hermite recovery}
\label{sec:cauchy}
\subsection{Deforming wire identifications instead of clause variables}
Within a packet, let $x,y\in\{0,\ldots,L-1\}$ be row and column positions.
Use the integer grid $S=\{-L,\ldots,L-1\}$ and set
\[
 \mathcal C(T)=\prod_{t\in S}(T-t),\qquad \alpha=2^{a+1},\qquad
 W_{xy}(T)=\frac{\mathcal C(T)}{\mathcal C'(0)(T-(y-x))}.
\]
The apparent pole cancels a numerator factor, so $W_{xy}$ is a polynomial
of degree $2L-1$. Between packets, set entries to zero; unused slots are
omitted. At zero, $W(0)=I$ on the real ports.

Let $A(Y)$ and $B(X)$ be the products of the row-side and column-side tensors,
respectively. With $l_0(z)=1-z$ and $l_1(z)=z$, define
\begin{equation}\label{eq:tensorF}
 F(T)=\sum_{X,Y\in\{0,1\}^M}A(Y)B(X)
       \prod_{\text{row ports }x}l_{Y_x}((W(T)X)_x).
\end{equation}
At $T=0$, the affine factors require the row and column bits to agree on
each original wire, so $F(0)$ is the original contraction. Its degree is at
most $M(2L-1)$. Formula \eqref{eq:tensorF} defines the polynomial; the
algorithm does not evaluate its double Boolean sum by brute force.

The grid is a complete residue system modulo $2^{a+1}$. For every node $t$,
the valuation of $\mathcal C'(t)$ equals that of $\mathcal C'(0)$.
Canceling these powers of two shows
\[
 W(t+\alpha z)\in\mathcal O[z],\qquad F(t+\alpha z)\in\mathcal A[z].
\]
More specifically, factors $t-q+\alpha z$ have known two-adic contents;
after cancellation only odd denominators need be inverted. This local
integrality is what makes modular evaluation legitimate, despite the rational
coefficients of the global Cauchy deformation.

\subsection{Multiplicity must vary with node valuation}
At nonzero node $t$, put $i=v_2(t)$ and $D_t=L/2^i$. Let
\[
 h=\lceil3M/u\rceil,\qquad q_t=h\max(0,u-i).
\]
There are $L$ odd nodes, so the total multiplicity exceeds $M(2L-1)$.
The source recovers $F(0)$ modulo $2^p$ from the first $q_t$ local
coefficients, using precision
\begin{equation}\label{eq:precision}
 p_t=p+8hD_t+8M(a+1).
\end{equation}
The taper $q_t$ is the key quantitative feature. Equal multiplicities at all
retained valuation levels would incur a larger precision loss at the last
level.

For each node define
\[
 E_t(T)=\prod_{q\in S\setminus\{0,t\}}
            \left(\frac{T-q}{t-q}\right)^{q_q}.
\]
Hermite reconstruction gives the identity
\begin{equation}\label{eq:taperidentity}
 F(0)=\sum_{t:q_t>0}E_t(0)
 \left[\frac{F(t+\alpha z)}{E_t(t+\alpha z)}\right]_{<q_t}
                  \bigg|_{z=-t/\alpha}.
\end{equation}
To verify it, replace zero by an indeterminate. Each summand has the desired
jet at its own node and vanishes to the required multiplicity at all other
nodes. The difference from $F$ has smaller degree than the product of all
those multiplicity factors and is divisible by that product, hence is zero.

The local denominator $E_t(t+\alpha z)$ has integral coefficients and constant
term one, so its formal reciprocal loses no precision. The source counts
grid points by their valuations and obtains
\begin{equation}\label{eq:valuationloss}
 v_2(E_t(0))=
 D_t\sum_{j=1}^{a-i}\frac{j}{2^j}\bigl(q(i+j)-q(i)\bigr)
 \ge-6D_th.
\end{equation}
Indeed, $q(i+j)-q(i)\ge-hj$ and $\sum_{j\ge1}j^2/2^j=6$.
Evaluation at $-t/\alpha$ loses at most another $4M(a+1)$ powers of two.
Thus \eqref{eq:precision} leaves enough precision for the final answer.
Rational arithmetic in this last combination has polynomial-length
numerators and denominators because all products and multiplicities are
polynomially bounded.

The recursive choice $a_d=a_{d-1}+2u_d$ ensures that inherited precision is
also affordable. At an active node $i<u_d$, one has
$D_t\ge2^{a_{d-1}+u_d}$. This additional factor dominates inherited
polylogarithmic overhead, yielding the effective bound
\begin{equation}\label{eq:effectiveprecision}
 p_t=O(MD_t/u_d).
\end{equation}
All precisions nevertheless remain polynomial in the root size $s$.
To obtain the local jets, the algorithm evaluates the entire local polynomial
at polynomially many residue-distinct extension-field lifts and interpolates
with unit denominators. Its evaluations are performed by the recursive
factorization described next.

\section{Valuation-weighted modes and the recursive contraction identity}
\label{sec:modes}
Fix a node $t$ with $i=v_2(t)<u$ and a local evaluation $T=t+\alpha\zeta$.
Group the vertices by coordinate $i$ of their labels. At $T=t$, a nonzero
Cauchy entry has $y-x=t$; adding an integer of valuation $i$ flips bit $i$.
The packet orientation convention therefore makes this entry join ports in
the same group. Only the perturbation away from the node creates crossing
interactions.

Partition a packet's columns into tiles
$y=y_0+2^{i+1}j$. A tile has a fixed source group. For a row crossing to
that tile, $v_2(t-y_0+x)\le i$. A geometric expansion gives
\begin{equation}\label{eq:modeexpansion}
 W_{xy}(T)\equiv
 \sum_{b+1<p_t}
 \frac{\mathcal C(T)}{\mathcal C'(0)(T-y_0+x)}
 \left(\frac{2^{i+1}}{T-y_0+x}\right)^b j^b
 \pmod{2^{p_t}}.
\end{equation}
After the known powers of two have been factored out, the inverses are unit
inverses. The coefficient at index $b$ is divisible by $2^{b+1}$.

For each packet, tile, and index $b$, introduce a mode $Z$ of weight
$w_Z=b+1$, whose eventual value is
\[
 L_Z(X)=\sum_{\text{real columns in the tile}}j^bX_y.
\]
Each mode has one source group and one target group. The row forms are sums
of a local part and these modes, with coefficients divisible by their
weight powers of two. Consequently a mode monomial $Z^\beta$ has coefficient
divisible by $2^{\sum_Zw_Z\beta_Z}$. It contributes nothing modulo $2^{p_t}$
if its total weight is at least $p_t$.

\subsection{A partition estimate counts all surviving interactions}
There are at most $6M/D_t$ modes of any one weight. For $\theta>0$, the number
of exponent tuples of total weight below $p_t$ is bounded by
\[
 \exp\left(\theta p_t+\frac{6M}{D_t}
                \sum_{j\ge1}-\log(1-e^{-\theta j})\right).
\]
Expanding the logarithm gives
\[
 \sum_{j\ge1}-\log(1-e^{-\theta j})
 =\sum_{h\ge1}\frac1{h(e^{\theta h}-1)}
 \le\frac1\theta\sum_{h\ge1}\frac1{h^2}.
\]
Choose $\theta=\sqrt{M/(D_tp_t)}$ and use
\eqref{eq:effectiveprecision}. The number of tuples is then
\begin{equation}\label{eq:modenumber}
 \exp(O(\sqrt{Mp_t/D_t}))=\exp(O(M/\sqrt u)).
\end{equation}
The tuples can be enumerated sequentially with polynomial-size counters.
A partial feasible tuple always has the all-zero continuation, so traversing
the enumeration tree has only a polynomial overhead per completed tuple.

\subsection{Fixing a tuple separates the groups}
For a fixed $\beta$, each target group extracts the coefficient of its
incoming modes from its affine row-factor product. Each source group
multiplies the corresponding powers $L_Z(X)^{\beta_Z}$. These operations
involve only bits from that group once $\beta$ is fixed.

To represent the resulting expression as a new bipartite network, put all
current vertices of one group on one side and new packet tensors on the
other. A packet tensor carries precisely the coefficient-extraction and
source-mode factors just described. Distributing the products and then
summing independent group bits gives the exact modular identity
\begin{equation}\label{eq:recursion}
 F(t+\alpha\zeta)
 =\sum_{\beta:\,\sum w_Z\beta_Z<p_t}
            \prod_{g=1}^{64}Z_{g,\beta}\pmod{2^{p_t}},
\end{equation}
where $Z_{g,\beta}$ is the contraction of that group's child network.
Terms not listed vanish by divisibility, not by an approximate truncation.

Each child has at most $6M/64=3M/32<M/8$ wires, using the balanced port loads.
A new packet vertex has at most $2L$ ports, so the degree invariant is
preserved. Its entry is computed by expanding a product of at most $L$
affine factors and extracting the specified coefficient, in
$\exp(O(L\log s))$ time and polynomial space. Since $L\le s^{1/1000}$,
this fits the allowed entry-access budget.

Descriptions are retained with their creation precision and used directly.
One does not recursively compose entry routines into an increasingly large
formula. This implementation detail prevents an uncounted description or
precision explosion along the recursion.

\section{Why the recursive costs sum to the square-root-logarithmic exponent}
\label{sec:tensorcost}
At depth $d$, $M\le Cs8^{-d}$ and
$u_d\asymp\lambda/2^d$. Thus \eqref{eq:modenumber} gives exponent
\[
 O\left(\frac{s}{\sqrt\lambda}
                 \left(\frac{\sqrt2}{8}\right)^d\right)
\]
for the branching work of a call, apart from polynomial factors for nodes,
evaluations, and groups. There are only $O(\log\lambda)$ active depths.
Multiplying the branching bounds adds their exponents, and the geometric
series converges. Polynomial overhead across these depths contributes only
$O(\lambda\log\lambda)$ to the logarithm of the cost, which is smaller than
$s/\sqrt\lambda$.

Leaves have at most $H\asymp s/\sqrt\lambda$ wires and are evaluated directly.
Their entry-access cost is $2^{O(s^{9/10})}$, also absorbed asymptotically by
$2^{O(s/\sqrt\lambda)}$. Correctness follows inductively from the exact
recursive identity, local modular evaluations, and tapered reconstruction.
The total bit-operation bound is therefore \eqref{eq:tensorbound}.

For space, use depth-first recursion. Retain one running sum over tuples,
one product over groups, the polynomially many local interpolation data,
and the current packet descriptions. Precision is polynomially bounded, and
the stack has only $O(\log\log s)$ levels. The computation therefore uses
polynomial space even though its total recursion tree is subexponentially
large.

The distinction from the rational clause-interpolation route is now visible.
There, a single global Taylor order produces a binomial subset-enumeration
cost. Here, modular divisibility controls a weighted mode count at each
recursive scale, and the tapered multiplicities reduce the precision enough
to make those costs sum geometrically. This is why the sharper exponent is
not obtained simply by deleting a logarithm from the earlier calculation.

\section{The arbitrary-density improvement via incidence-subset encoding}
\label{sec:tensordense}
The tensor theorem must next be applied to sparse polynomial parity systems
without creating one wire per incidence, since there can be
$n\operatorname{polylog}(n)$ incidences. The source uses the global degree bound
\eqref{eq:totaldegree} to encode selected incidences more economically
\cite{quant-reduction}.

After branching on heavy variables as before, introduce a separate variable
$z_{iv}$ for each test--variable incidence. Multilinearize each lifted test
factor in its own incidence variables. The factors now use disjoint
indeterminate sets, and their truncated product is multilinear of total
incidence degree at most
\[
 K=k(a+2n)=O(n).
\]
The actual Boolean variable $x_v$ is substituted for all of its incidence
copies only after this coefficient representation has been established.

Partition the tests and variables separately into
$J=\lceil n^{1/4}\rceil$ nearly equal groups. Let $N_{gh}$ be the number of
incidences between test group $g$ and variable group $h$, with total
$N\le n(\log n)^{O(1)}$. For a monomial, write $d_{gh}$ for its selected
incidence count in that pair. Enumerate arrays with
$0\le d_{gh}\le N_{gh}$ and $\sum d_{gh}\le K$. There are at most
$(K+1)^{J^2}=2^{O(\sqrt n\log n)}$ such arrays.

For a fixed array, the subset in pair $(g,h)$ has
$\binom{N_{gh}}{d_{gh}}$ possibilities. Encode its lexicographic index with
$\lceil\log_2\binom{N_{gh}}{d_{gh}}\rceil$ Boolean wires. Invalid bit strings
receive tensor value zero. Subsets are decoded by binomial-coefficient counts
of completions, in polynomial time. The total wire count is bounded by
\begin{align}\label{eq:wirecount}
 M&\le J^2+\log_2\prod_{g,h}\binom{N_{gh}}{d_{gh}}\\
  &\le J^2+\log_2\binom{N+K}{K}
   =O(n\log\log n).\nonumber
\end{align}
The last estimate uses $K=\Theta(n)$ and $N/n$ polylogarithmic.
This entropy bound is the source of the single remaining $\log\log n$
factor in the sharper arbitrary-density exponent.

At a test-group vertex, the tensor entry is the product of the coefficients
of the selected incidence monomials in that group's test factors. A selected
coefficient is computed by subset inversion:
\[
 [z_S]\overline B_i(z,V)
   =\sum_{T\subseteq S}(-1)^{|S|-|T|}B_i(\one_T,V).
\]
Each test scope has polylogarithmic size, so this entry procedure costs only
$2^{(\log n)^{O(1)}}$ time and polynomial space.
At a variable-group vertex, a variable selected at least once must equal one
when summing its monomial over the Boolean cube; an unselected variable is
free. The tensor value is therefore $2$ to the number of unselected variables,
computed modulo $2^{n+1}$.

Every incidence monomial occurs in exactly one degree array and has one valid
subset index per group pair. The product of the test-side entries is its
coefficient, and the product of the variable-side entries is its Boolean
monomial sum. Summing the resulting contractions therefore gives the required
truncated polynomial sum exactly. Apply $\mathcal S_\ell$ only at the end.

Each group has at most $n^{3/4}(\log n)^{O(1)}$ incident ports, within the
contraction theorem's degree allowance for large $n$ when
$s=n\lceil\log_2\log_2n\rceil$. Its precision and marker truncation are
$n+1=O(s)$. Equations \eqref{eq:tensorbound} and \eqref{eq:wirecount} thus give
\[
 n^{O_{k,c}(1)}2^{O_{k,c}(n\log\log n/\sqrt{\log n})}
\]
for the sparse polynomial counter. Heavy-variable branches and the degree-array
enumeration have smaller exponents and are processed sequentially.

Finally use disjoint sparsification with
$\sigma\asymp1/\sqrt{\log n}$, and the same deterministic parity-test selection,
now calling this faster counter. The leaf count and candidate-enumeration
costs are absorbed in \eqref{eq:sharp}. A single algorithm can achieve both
claimed regimes: use the direct clause network when
$m\le n\lceil\log_2\log_2n\rceil$ and the parity-system route otherwise.
For every fixed density $D$, all sufficiently large inputs with $m\le Dn$
choose the first branch, giving \eqref{eq:density}.

\section{The explicit ETH contradiction and the canonical-input issue}
\label{sec:eth-consequence}
Take $k=3$. The reported algorithm computes $\#F$ exactly; comparing its
binary output with zero decides satisfiability. Even the weaker exponent
\eqref{eq:rho} is enough because
\[
 \frac{\rho(n)}n=\frac{(\log\log n)^2}{\sqrt{\log n}}\longrightarrow0.
\]
For any proposed constant $c>0$ in \eqref{eq:eth}, a fixed exponent constant
$A$ satisfies $A\rho(n)\le cn/2$ for all sufficiently large $n$.
The finitely many smaller inputs can be absorbed into the overall constant.
Thus one obtains an algorithm within the rate that \eqref{eq:eth} forbids.
The source's stronger uniform statement chooses one machine for fixed $k$
before all input sizes and formulas; it is not advice selected separately
for every input.

The published formal theorem uses a canonical encoding. Variables lie in
$\operatorname{Fin}(n)$, clauses and formulas are finite sets, and literals,
clauses, and formulas are sorted and framed. Natural-number components use
unary coding. Therefore the theorem's polynomial factor is measured in this
canonical length $L$, not in an unspecified compact representation.

For an ordinary explicit 3-CNF decision input, rename its occurring variables
consecutively. Their number is at most the input length, and each unary index
then has polynomial length in that input. Sorting, removing duplicate literals
and clauses, and producing the canonical representation are polynomial-time
operations. The polynomial factor in $L$ is consequently a polynomial
input-reading factor in the usual representation. This supplies the needed
bridge to the catalogue's decision formulation.

For fixed width there are only
$\sum_{j=0}^k\binom{2n}{j}=O_k(n^k)$ distinct normalized clauses. Thus the
canonical representation also cannot hide a positive exponential rate solely
in its duplicate-free clause list. Exact counting must still retain unused
declared variables or restore a factor $2^u$. A declaration of an astronomically
large number of unused variables in a compact binary field requires special
care about output length: writing the count itself needs up to $n+1$ bits.
The formal encoding and the manuscript's polynomial-in-$n+m$ statement include
that cost. No stronger succinct-input counting claim is inferred here.

\section{Consequences that follow, and claims that do not}
\label{sec:consequences}
\subsection{Witnesses and exact sampling}
A satisfying assignment can be found by fixing one bit at a time and counting
the two restrictions. If the current count is positive, at least one branch
has positive count. This needs only polynomially many calls.
To sample uniformly, choose each next bit with probability equal to its
conditional count divided by the current count. For a final satisfying
assignment the conditional ratios telescope to $1/\#F$. Exact integer
sampling implements these rational probabilities with polynomial expected
overhead. This is a randomized exact sampler, not a deterministic procedure
that somehow produces randomness.

\subsection{Examples of parameter-controlled counting reductions}
Independent sets of a graph correspond bijectively to satisfying assignments
of $\bigwedge_{\{u,v\}\in E}(\neg x_u\vee\neg x_v)$, using one variable per
vertex. For each fixed number of colors $q$, one-hot color variables give a
count-preserving fixed-width encoding with $q$ variables per vertex. These
reductions therefore give subexponential counting in the number of vertices,
up to the relevant polynomial input factors.

For perfect matchings, one Boolean variable per edge and a linear-size
checking circuit give a unique-extension width-three encoding with
$O(v+e)$ variables. The resulting subexponential parameter is $v+e$, not
automatically $v$ for dense graphs. This distinction illustrates why the
size of a reduction must be tracked rather than suppressed in the phrase
``NP-complete reduction'' \cite{paper}.

\subsection{What ETH's failure would and would not imply}
A subexponential algorithm need not be polynomial time. For instance,
$2^{\sqrt n}$ is subexponential and superpolynomial. Thus the reported theorem
does not settle $\mathsf P$ versus $\mathsf{NP}$. Its constants can also be
large, and none of the proofs supplies a practical SAT benchmark or a small
crossover input size.

The fixed-width theorem contradicts the usual fixed-$k$ strong exponential-time
hypothesis as well: every fixed width now has vanishing exponential rate.
It does not assert a uniform subexponential bound for unrestricted-width CNF
with $k$ growing with $n$. Nor does it automatically give subexponential
algorithms for every problem in NP. A polynomial reduction can enlarge an
$n$-parameter instance to $n^c$ SAT variables, and $2^{o(n^c)}$ need not be
$2^{o(n)}$.

An ETH-conditional lower-bound theorem remains a valid implication if ETH is
false, but its antecedent can no longer establish an unconditional lower bound.
To obtain a fast algorithm for its target problem one still needs a suitable
reduction in the algorithmic direction with a controlled parameter blowup.
The source contains further applications to approximation, local maps, and
circuit lower bounds; those are downstream results, not prerequisites for
the CS-20 resolution, and their full proofs are not reproduced in this report.

\section{Formal scope, source dependencies, and remaining verification}
\label{sec:formal}
The inspected comparator target quantifies a fixed deterministic primitive
machine for each $k\ge2$, then every $n$ and every width-at-most-$k$ formula.
A basic machine step performs a bounded tape action, rather than an arbitrary
unit-cost rational calculation. The target supplies a halting trace with the
correct binary count, the satisfiability bit, the documented runtime, and a
polynomial bound on total visited tape cells at every execution prefix
\cite{challenge}.

The runtime scale in that statement uses shifted logarithms,
\[
 n\frac{[\log\log(n+16)]^2}{\sqrt{\log(n+16)}},
\]
which is asymptotically the same as \eqref{eq:rho} and well-defined for all
natural-number inputs. The formal scope explicitly distinguishes this bound
from the sharper manuscript exponents \cite{scope}. This report likewise
does not label the tensor-contraction runtime as independently machine-checked
merely because a weaker sufficient SAT theorem is documented.

The reference file under \texttt{ComparatorChallenges} contains the target
statement and its proof placeholder. It is not the solution proof. Its JSON
configuration routes to the separate module
\texttt{OAI.Algorithms.SubexponentialSAT.Main} and names the theorem
\begin{center}\small
\texttt{OAI.Algorithms.SubexponentialSAT.SqrtLogSpace.}\par
\texttt{exact\_count\_and\_sat}.
\end{center}
The permitted foundational axioms are \texttt{propext}, \texttt{Quot.sound},
and \texttt{Classical.choice} \cite{config,solution}.
Inspecting these files is not the same as successfully executing their
verification procedure. No complete rebuild, comparator run, or transitive
axiom audit was performed here.

The mathematical dependencies that most deserve an independent check are the
balanced-incidence embedding and denominator bounds in the rational route;
the global marker-degree truncation and noncircular parity-test selection;
and, in the sharper route, local integrality, the valuation loss in tapered
Hermite recovery, the child-network identity, and the immutable precision
convention for implicit tensors. Each is necessary to rule out a hidden
exponential cost or an invalid algebraic substitution.

The earlier catalogue contribution explored separator-based strategies and
correctly distinguished them from a uniform algorithm. The source method
reported here does not assume those separators exist: its separations occur
at algebraically deformed inputs or within valuation-indexed recursive
contractions. This is the substantive difference from an unproved graph
partition hypothesis \cite{catalogue}.

\section{Completion Estimate}
\noindent\textbf{COMPLETION: 0\%}
This report claims no new original solution or independently verified proof. The estimate concerns this report's original research contribution, not the extent of the results attributed to OpenAI.
\section{Final status and preparation}
Unresolved as an original research attempt: literature exposition only. OpenAI's reported full result remains attributed to the cited source; no source-catalogue or community-review verdict is changed by this report.
The document is dated 26 October 2026. The underlying GPT 6 Astra Pro release/build and reasoning setting are unknown. The supplied source-inspection statements describe the preparation of the report; repository integration checks identity, formatting, compilation and rendering, rather than certifying the underlying proofs. Original generation prompts and browsing logs are unavailable; use of the repository's initial and final research prompts during original preparation is not documented. No independent human audit, Lean execution or new performance benchmark is asserted.


\section{Attribution and provenance}
\paragraph{Problem handle.} \texttt{prb\_a1b2b1f67d0e489da50d97b732d8291d}
(CS-20).
\paragraph{Submission claim.} Literature report.
\paragraph{Underlying reported result.} Full disproof of ETH, attributed to
OpenAI, through deterministic exact counting for each fixed clause width.
\paragraph{Exact scope.} Detailed exposition of the bounded-occurrence
interpolation core, disjoint general reduction, sparse-polynomial counter,
deterministic parity tests, arithmetic and space accounting, and the sharper
recursive tensor method. It derives the ETH contradiction and checks ordinary
input compatibility. It distinguishes the documented formal runtime from the
sharper manuscript runtimes. No original discovery, full independent audit,
Lean execution, or practical performance claim is made.
\paragraph{Who did the work.} LLM only (people only prompted or submitted).
\paragraph{Model used in this preparation.} OpenAI GPT 6 Astra Pro; exact version:
unknown; thinking setting: unknown; actual time or compute budget: unknown.
Role: source inspection, mathematical exposition, scope comparison, and TeX
preparation. No second model or independent human proof reviewer was used.
\paragraph{People and roles.} The submitter prompted preparation and expansion. No human mathematical authorship or independent proof-verification role is asserted.
\paragraph{External research provenance.} OpenAI is the author named on the
source manuscript. The original research's internal model identities, settings,
budgets, and individual roles are not established by this report and are
recorded as unknown. They are external source provenance, not additional models
invoked during this preparation.
Linked third-party artifacts retain their original terms. Source manuscripts, images and Lean code are not bundled.

\begin{thebibliography}{99}
\bibitem{catalogue} \emph{we-solve}, CS-20,
\emph{The Exponential Time Hypothesis for 3-SAT}.
Used: exact ETH quantifiers, input-size convention, stable handle, and the
scope of the earlier separator exploration.
\url{https://we-solve.github.io/problems/the-exponential-time-hypothesis-for-3-sat/}.
\bibitem{paper} OpenAI, \emph{The Exponential Time Hypothesis is False},
October 17, revised October 25, 2026, pinned commit in Section~\ref{sec:target}.
\href{https://github.com/openai/math/blob/7ef5c0cc5ca4b0e072b7a050fca09ab24bc6fa3a/preprints/The-Exponential-Time-Hypothesis-is-False-October-25-2026/build/main_full_version.tex}{Main manuscript source}.
Used: theorem and technique overview; measurement separation; Boolean Taylor
recovery; swap-matrix construction; incidence balancing; parity lift and its
total-degree bound; sparse-polynomial counting; deterministic parity tests;
sparsification appendix; bit complexity and polynomial space; selected
counting consequences. The source was read for this exposition, not fully
independently verified.
\bibitem{quant-reduction} OpenAI, same manuscript,
\href{https://github.com/openai/math/blob/7ef5c0cc5ca4b0e072b7a050fca09ab24bc6fa3a/preprints/The-Exponential-Time-Hypothesis-is-False-October-25-2026/build/quantitative-reduction.tex}{Sharper exact-counting bounds and tensor encodings}.
Used: contraction interface, fixed-density clause network, incidence-subset
encoding, dense-case reduction, and the single algorithm attaining both regimes.
\bibitem{quant-contraction} OpenAI, same manuscript,
\href{https://github.com/openai/math/blob/7ef5c0cc5ca4b0e072b7a050fca09ab24bc6fa3a/preprints/The-Exponential-Time-Hypothesis-is-False-October-25-2026/build/quantitative-contraction.tex}{Proof of the Boolean tensor-contraction theorem}.
Used: coefficient arithmetic, packet packing, Cauchy deformation, local
integrality, tapered Hermite recovery, weighted modes, child networks,
precision invariants, and recursive time/space bounds.
\bibitem{scope} OpenAI,
\href{https://github.com/openai/math/blob/7ef5c0cc5ca4b0e072b7a050fca09ab24bc6fa3a/lean/docs/378.md}{Formalization scope 378}.
Used: exact-counting and satisfiability scope, canonical input length,
documented weaker runtime, and the distinction from the manuscript bounds.
\bibitem{challenge} OpenAI,
\href{https://github.com/openai/math/blob/7ef5c0cc5ca4b0e072b7a050fca09ab24bc6fa3a/lean/ComparatorChallenges/SubexponentialSAT.lean}{Subexponential-SAT comparator reference}.
Used: finite variable set, canonical encoding, primitive-machine semantics,
quantifier order, binary output, runtime, and all-prefix space statement.
This is the reference specification, not its solution proof.
\bibitem{config} OpenAI,
\href{https://github.com/openai/math/blob/7ef5c0cc5ca4b0e072b7a050fca09ab24bc6fa3a/lean/ComparatorChallenges/SubexponentialSAT.json}{Subexponential-SAT comparator configuration}.
Used: theorem identification, solution routing, and permitted axioms.
\bibitem{solution} OpenAI,
\href{https://github.com/openai/math/blob/7ef5c0cc5ca4b0e072b7a050fca09ab24bc6fa3a/lean/OAI/Algorithms/SubexponentialSAT/Main.lean}{Subexponential-SAT main solution module}.
Used: imports and identification of the exact-counting polynomial-space result.
No complete dependency-closure check or successful local execution is claimed.
\end{thebibliography}
\end{document}
